\documentclass[10pt,a4paper]{amsart}
\usepackage[margin=19mm]{geometry}
\usepackage{amsmath,amssymb,mathtools}
\usepackage[scr=rsfs,cal=euler]{mathalfa}
\usepackage{xfrac}
\usepackage[unicode,hidelinks,bookmarks=false]{hyperref}
\newcommand{\ie}{i.e.\ }
\newcommand{\cf}{cf.\ }
\newcommand{\resp}{respectively\ }
\newcommand{\Subsection}{\subsection}
\providecommand{\nobreakdash}{\nobreak\hbox{-}}
\setlength{\emergencystretch}{2em}
\multlinegap0pt
\theoremstyle{plain}
\newtheorem{thm}{Theorem}
\newtheorem{prop}{Proposition}
\newtheorem{lem}{Lemma}
\newtheorem{cor}{Corollary}
\theoremstyle{definition}
\newtheorem{defn}{Definition}
\theoremstyle{remark}
\newtheorem{rem}{Remark}
\title[3+1d symmetry enriched topological order]{The Classification of 3+1d Symmetry Enriched Topological Order: two-categorical (de-)equivariantization}
\author{Thibault D. Décoppet, Matthew Yu}
\date{Source: arXiv:2509.10603v3 (CC BY 4.0)}
\begin{document}
\maketitle
\noindent\emph{Source and licence.} The Classification of 3+1d Symmetry Enriched Topological Order. Version arXiv:2509.10603v3 (CC BY 4.0), distributed under the Creative Commons Attribution 4.0 International licence, \url{http://creativecommons.org/licenses/by/4.0/}, as recorded in the arXiv licence metadata for this exact version and shown on the abstract page \url{https://arxiv.org/abs/2509.10603v3}. Full licence notice: licenses/arxiv-2509.10603-CC-BY-4.0.txt.\par\smallskip

\noindent\textit{Editorial scope.} Counted unit: the article's Theorem with source label \texttt{thm:deequiv} (source0.tex lines 745--747), the two-categorical (de-)equivariantization theorem identifying braided fusion 2-categories containing the two-representations of the group fully faithfully with the corresponding crossed braided multifusion 2-categories; delivered together with the article's complete original proof of it (source0.tex lines 748--758, one \texttt{proof} environment of 11 lines, adjacent to the statement, which the article presents as the categorification of the classical argument for fusion categories). No mathematics is written, repaired or re-derived: statement and proof are the article's own text line for line, in the article's own notation (braided fusion 2-categories, crossed braided multifusion 2-categories, rigid braided algebras, Drinfeld centres and higher Morita equivalence), and no retained line is substituted, re-flowed or omitted. No further part of the article is required as context and none is retained: the counted proof uses only the objects introduced in its own statement together with the standard Morita-invariance of Drinfeld centres that it names, so no cross-reference and no citation occurs in any retained line; consequently the wrapper carries no bibliography and no bibliography entry. Nothing else of the article is retained: its classification results, its other propositions and its diagrams are omitted. Editorial formatting only: an AMS \texttt{amsart} preamble that restates the environments the retained lines use, namely the article's declarations of Theorem, Proposition, Lemma, Corollary, Definition and Remark, and the macros that the retained lines need. The article's own source declares the class \texttt{amsart} as well, and its \texttt{quiver} style file is neither loaded nor redistributed here; no class or style file is shipped with this package. Numbering differs from the article, because only one environment is retained and the wrapper keeps a continuous counter: the counted theorem prints as Theorem 1. No picture environment and no graphical inclusion occurs in any retained line, and the package delivers no image asset. One bundled source file, \texttt{sources/184985/source0.tex}, accompanies the delivered item as the exact source of every retained span. Every retained span is recorded in \texttt{delivery-evidence.json} with method \texttt{exact\_tex}.
\begin{thm}{(2-categorical (de-)equivariantization)}\label{thm:deequiv}
There is an equivalence of 3-categories between braided fusion 2-categories containing $\mathbf{2Rep}(G)$ fully faithfully and $G$-crossed braided fusion 2-categories.
\end{thm}

\begin{proof}
By definition, we have that $G$-crossed braided multifusion 2-categories are rigid braided algebras in $\mathcal{Z}(\mathbf{3Vect}_G)$. On the other hand, we have $\mathbf{Mod}(\mathbf{2Rep}(G))\simeq \mathbf{3Rep}(G)$ as (symmetric) fusion 3-categories as the right hand-side is connected. Moreover, rigid braided algebras in $\mathcal{Z}(\mathbf{Mod}(\mathbf{2Rep}(G)))$ are braided multifusion 2-categories equipped with a braided 2-functor from $\mathbf{2Rep}(G)$. Now, we have that the fusion 3-categories $\mathbf{3Vect}_G$ and $\mathbf{3Rep}(G)$ are Morita equivalent via $\mathbf{3Vect}$. It follows that there is an equivalence of braided fusion 3-categories
\begin{equation}\label{eq:Z3VGZ3RG}
\mathcal{Z}(\mathbf{3Vect}_G)\simeq\mathcal{Z}(\mathbf{3Rep}(G)).
\end{equation}
This is a consequence of the fact that Morita equivalent fusion higher categories have equivalent Drinfeld centers.
This establishes that the 3-category of $G$-crossed braided multifusion 2-categories is equivalent to the 3-category of braided multifusion 2-categories equipped with a braided 2-functor from $\mathbf{2Rep}(G)$.
Namely, braided rigid algebras are preserved by arbitrary braided monoidal functors.

In order to conclude the proof, it remains to show that this correspondence sends $G$-crossed braided fusion 2-categories to braided fusion 2-categories equipped with a fully faithful braided 2-functor from $\mathbf{2Rep}(G)$. We have recalled above that an algebra in a higher fusion category is called strongly connected if its unit 1-morphism is the inclusion of a summand. It is clear that strongly connected rigid algebras in $\mathcal{Z}(\mathbf{3Vect}_G)$ are $G$-crossed braided fusion 2-categories. On the other hand, strongly connected rigid algebras in $\mathcal{Z}(\mathbf{Mod}(\mathbf{2Rep}(G)))$ are braided fusion 2-categories equipped with a fully faithful braided 2-functor from $\mathbf{2Rep}(G)$. The proof is completed by appealing to Equation \eqref{eq:Z3VGZ3RG}.
\end{proof}

\end{document}
